\section{Conclusion}
\label{sec:conclusion}

We introduce \method,
a framework that is capable of jointly reconstructing dense geometry and motion from a monocular video.
By defining both in a unified world coordinate system
and designing a novel 4D VAE that encodes them into a shared latent space,
we achieve state-of-the-art performance even without any post-processing. 
Notably,
we show that it is not necessary to strictly align the 4D latent distribution with the original SVD latent distribution.
In fact,
our relaxed alignment and VAE retraining strategy not only preserves but also improves the diffusion model's generalization,
offering broader insight into adapting diffusion priors to new modalities.

\paragraph{Limitations}
Currently,
we focus solely on dense geometry and motion reconstruction,
but prior work has shown that incorporating multiple geometric modalities can substantially improve the prediction of 3D attributes, including camera parameters, point maps, depth maps, point tracks, and novel views.
Thus, exploring multimodal integration is a promising direction for future work.

\section*{Acknowledgments}
% Ruijie Zhu completed this work during his internship at Tencent ARC Lab. We thank Tian-Xing Xu for providing the codebase of GeometryCrafter.
Xiaoguang Han is supported by Guangdong Provincial Outstanding Youth Fund with No. 2023B1515020055.
Chuanxia Zheng is supported by NTU SUG-NAP and the National Research Foundation, Singapore, under its NRF Fellowship Award NRF-NRFF17-2025-0009.
\\